\documentclass[t,usenames,dvipsnames]{beamer} % [t] pomeni poravnavo na vrh slida

    % \State Randomly choose rewriting rule 
    % \State choose rewriting rule Randomly 
    % relation closure more verbose on slides

% \usepackage{etex} % vključi ta paket, če ti javi napako, da imaš naloženih preveč paketov.
% \usefonttheme[stillsansserifsmall,stillsansseriflarge]{serif}

% \usepackage{amsfonts}
% \usepackage[mathscr]{eucal}
% \usepackage[mathscr]
\usepackage{mathrsfs}


% standardni paketi
\usepackage[slovene]{babel}
\usepackage[T1]{fontenc}
\usepackage[utf8]{inputenc}
\usepackage{amsmath,amssymb,amsfonts,amsthm} % matematični paketi
% \usepackage{amssymb}
% \usepackage{amsmath}
\usepackage{diagrams}
\usepackage{array}
\usepackage{multirow}
\usepackage{qtree}
	\usepackage{forest}
\forestset{
dg edges/.style={for tree={parent anchor=south, child anchor=north,align=center,base=bottom,where n children=0{tier=word,edge=dotted,calign with current edge}{}}},
dgii edges/.style={for tree={parent anchor=south, child anchor=north,base=bottom}},
}


\DeclareMathOperator{\Coverage}{Coverage} 

% % down 4.4.2017
\newcommand{\R}{\mathbb R}
% \newcommand{\N}{\mathbb N}
% \newcommand{\Z}{\mathbb Z}
\newcommand{\C}{\mathbb C}
\newcommand{\Q}{\mathbb Q}
% a51303
\definecolor{softyellow}{rgb}{0.98,0.98,0.75}
% \definecolor{dunered}{HTML}{a51303}
\definecolor{dunered}{HTML}{eb4702}
\definecolor{dunered}{HTML}{912010}
\definecolor{dunered}{HTML}{bd1a03}

\setbeamercolor{loweryel}{fg=black,bg=softyellow}
% % up 4.4.2017
\newcommand{\e}{\boldsymbol{e}}
\newcommand{\X}{\mathscr{X}}
\newcommand{\I}{\mathscr{I}}
\newcommand{\rhoo}{\boldsymbol{\rho}}
\newcommand{\q}{\boldsymbol{q}}
\newcommand{\1}{\boldsymbol{1}}
\newcommand{\0}{\boldsymbol{0}}
\newcommand{\N}{\mathbb{N}}
\newcommand{\Nom}{\mathbb{N}_0^m}
\newcommand{\Z}{\boldsymbol{Z}}
\newcommand{\Zn}{\boldsymbol{Z}_n}
\newcommand{\M}{\boldsymbol{M}}
\newcommand{\re}{\boldsymbol{r}}
% \newcommand{\r}{\boldsymbol{r}}
\renewcommand{\r}{\boldsymbol{r}}
\newcommand{\s}{\boldsymbol{s}}
% glej $$ \r $$ 
% \newcommand{\e}{\boldsymbol{e}}
\renewcommand{\e}{\boldsymbol{e}}

% \newtheorem{lemma}{Lemma}
% \newtheorem{definition}{Definition}
% \newtheorem{theorem}{Theorem}



%\setbeamercovered{invisible} %default
\setbeamercovered{transparent}
%\setbeamercovered{dynamic}


%\usepackage[dvipsnames]{color}

\usepackage[normalem]{ulem} % za strikeout (prečrtat besedo)
% na primer : \sout{Hello World}

% podatki
\title{Discovery of exact equations via \\ computing Gr\"{o}bner basis }
\author{Bo\v{s}tjan Gec}
\institute{ mentor: prof. dr. Ljup\v{c}o Todorovski \\
            comentor: prof. dr. Sa\v{s}o D\v{z}eroski}

% tvoj izbran stil predstavitve
%\usetheme{Singapore}
\usetheme{Luebeck}
%\usecolortheme{crane}
%\usecolortheme{red}

\usepackage{color}
\setbeamercolor{structure}{fg=Bittersweet}
\setbeamercolor{structure}{fg=ForestGreen}


% \usepackage{algorithmic}
\usepackage{algpseudocode}  % za psevdokodo
\usepackage{algorithm}
% \floatname{algorithm}{Algoritem}
% \renewcommand{\listalgorithmname}{Kazalo algoritmov}
\usepackage{multirow}

\usepackage{hyperref}

\usepackage{booktabs}%



\newcommand{\MB}{\textsl{MoadeeB}}
\newcommand{\MoBu}{M\"oller-Buchberger}
\newcommand{\MoBua}{M\"oller-Buchberger algorithm}
\newcommand{\Bua}{Buchberger algorithm}
\newcommand{\DP}{Diofantos}
\newcommand{\linrec}{\textsl{linrec}}  % data set
\newcommand{\core}{\textsl{core}}  % data set
\newcommand{\GB}{Gr\"obner basis}


\begin{document}


\begin{frame}
  \maketitle
%  \fillblack{0.33}
\end{frame}


% \section{Equation discovery}
% \section{Verjetnostne gramatike}
% \section{Celo"stevilska zaporedja}

\begin{frame}{\GB\ in action}

    \begin{itemize}[<+->]
        \item Variables $V = \{a_n, a_{n-1}, a_{n-2}\}$.
        \item Polynomials over $\Q$, i.e. in $\mathbb{Q}[a_n, a_{n-1}, a_{n-2}]$.
        \item E.g. polynomial $$ p(x,y,z) = x^2 + y z $$ ... looks like  ...
            $$f = a_n^2 + a_{n-1}a_{n-2}$$ in $\mathbb{Q}[a_n, a_{n-1}, a_{n-2}]$.
    \end{itemize}

\end{frame}


\begin{frame}{Inductively build the vanishing ideal}

    Vanishing ideal:
$$ I(X) := \{ f \in K[X_1, ..., X_n] \mid f(x_1, ..., x_n) = 0 \text{ for all } (x_1, ..., x_n) \in X \}\enspace $$ 

    \begin{block}{$n=1$}
        one data point ... is a variety itself \\ \uncover<2>{
        ... corresponding ideal $I_1 = \langle a_n - 1, a_{n-1} - 1, a_{n-2} \rangle$.
    }
    \end{block}

    \begin{block}{$n \to n+1$}
        Union of varieties ... product of ideal bases \uncover<2>{($I_1 \cdot I_2$) }
    \end{block}
\end{frame}

\begin{frame}{Inductively build the vanishing ideal}

    \begin{block}{$n=1$}
        one data point ... is a variety 
        $V_1 := \{ (1, 1, 0) \} $

        \begin{center}
        $\updownarrow$ corresponds to
        \end{center}

        an ideal ... $I_1 = \langle a_n - 1, a_{n-1} - 1, a_{n-2} \rangle$.
        \pause

$$ I(X) := \{ f \in K[X_1, ..., X_n] \mid f(x_1, ..., x_n) = 0 \text{ for all } (x_1, ..., x_n) \in X \}\enspace,  $$ 

        \pause

         i.e. \begin{align*} I(V_1) = \{ p \in \Q[a_n, a_{n-1}, a_{n-2}] \mid p(a_n, a_{n-1}, a_{n-2}) = 0  \text{ if }  \\  (a_n, a_{n-1}, a_{n-2}) = (1,1,0) \ \} \end{align*}

        \pause

         i.e. \( I(V_1) = \{ p \in \Q[a_n, a_{n-1}, a_{n-2}] \mid p(1,1,0) = 0 \} \)
    \end{block}

    \begin{block}{$n \to n+1$}

    \end{block}


\end{frame}

\begin{frame}{Inductively build the vanishing ideal}

    \begin{block}{$n=1$}
        one data point ... is a variety 
        $ \{ (1, 1, 0) \} $

        \begin{center}
        $\updownarrow$ 
        \end{center}

        an ideal ... $I_1 = \langle a_n - 1, a_{n-1} - 1, a_{n-2} \rangle$.

        \pause

        \begin{itemize}[<+->]
            \item
        Fun fact:
        $$a_n - a_{n-1} - a_{n-2} \in I_1$$
            \item
        Def. of ideal:
for any polynomials \(p, q \in I\) and \(r \in \mathbb{Q}[a_n, a_{n-1}, a_{n-2}]\), the polynomials $0$, \(p+q\) and \(r \cdot p\) are also in $I$.
        \end{itemize}

    \end{block}


\end{frame}


\begin{frame}{Inductively build the vanishing ideal}

    \begin{block}{$n=1$}
        one data point ... an ideal $I_1 = \langle a_n - 1, a_{n-1} - 1, a_{n-2} \rangle$.
    \end{block}

    \begin{block}{$n \to n+1$}

        \begin{itemize}
            \item Union of varieties ... product of ideal bases ($I_1 \cdot I_2$) \\

        \pause

        \item
        Basis of ideal of union of two varieties ... obtained by pairwise multiplications of the polynomials in the bases of ideals of the individual varieties. 
        \pause

        \item
            $ I( V_1 \cup V_2 ) = I(V_1) \cdot I(V_2) = $ \\ \pause
            e.g. $  = \langle f_1, f_2 \rangle \cdot \langle g_1, g_2 \rangle = $ \\ \pause
            $  = \langle f_1\cdot g_1, f_1\cdot g_2, f_2\cdot g_1, f_2\cdot g_2 \rangle = $ \\ \pause
                $I(\{(1,1,0)\} \cup \{(2,1,1)\}) = I(\{(1,1,0)\}) \cdot I(\{(2,1,1)\}) = $ \\ \pause
               $ = \langle a_n - 1, a_{n-1} - 1, a_{n-2} \rangle \cdot \langle a_n - 2, a_{n-1} - 1, a_{n-2}-1 \rangle = $ 


        \end{itemize}


    \end{block}


\end{frame}


\begin{frame}{Inductively build the vanishing ideal}

    \begin{block}{$n=1$}
        one data point ... an ideal $I_1 = \langle a_n - 1, a_{n-1} - 1, a_{n-2} \rangle$.
    \end{block}

    \begin{block}{$n \to n+1$}

        \begin{itemize}
        \item
                $ I( V_1 \cup V_2 ) = I(V_1) \cdot I(V_2) = $ \\ 
        \uncover<3>{ e.g. $  = \langle f_1, f_2 \rangle \cdot \langle g_1, g_2 \rangle = $ \\ 
        $  = \langle f_1\cdot g_1, f_1\cdot g_2, f_2\cdot g_1, f_2\cdot g_2 \rangle = $ \\ 
                }
        
        \item
        \item
                $I(\{(1,1,0)\} \cup \{(2,1,1)\}) = I(\{(1,1,0)\}) \cdot I(\{(2,1,1)\}) = $ \\  
               $ = \langle a_n - 1, a_{n-1} - 1, a_{n-2} \rangle \cdot \langle a_n - 2, a_{n-1} - 1, a_{n-2}-1 \rangle = $   \pause
    \begin{align*}
     = & \langle (a_n -1)(a_n -2), (a_n -1)(a_{n-1} -1), (a_n -1)(a_{n -2}-1), \\ 
    & (a_{n-1} -1)(a_n -2),  (a_{n-1} -1)(a_{n-1} -1), (a_{n-1} -1)(a_{n -2}-1), \\
    & (a_{n -2})(a_n -2), (a_{n -2})(a_{n-1} -1), (a_{n -2})(a_{n -2}-1) \rangle = \\
    \end{align*}

        \item Literature~\cite[p.~24, Section~2.1, Chapter~1]{Shafarevich}.

        \end{itemize}

    \end{block}
\end{frame}

\begin{frame}{Inductively build the vanishing ideal}

    \begin{block}{$n=1$}
        one data point ... an ideal $I_1 = \langle a_n - 1, a_{n-1} - 1, a_{n-2} \rangle$.
    \end{block}

    \begin{block}{$n \to n+1$ }
        \begin{itemize}
        \item
                $I(\{(1,1,0)\} \cup \{(2,1,1)\}) = I(\{(1,1,0)\}) \cdot I(\{(2,1,1)\}) = $ \\
               $ = \langle a_n - 1, a_{n-1} - 1, a_{n-2} \rangle \cdot \langle a_n - 2, a_{n-1} - 1, a_{n-2}-1 \rangle = $  
    \begin{align*}
    = & \langle a_n^2 -3a_n +2,  a_n a_{n-1} -a_n -a_{n-1} +1,  a_n a_{n-2} -a_n -a_{n-2} +1, \\
    & a_n a_{n-1} -a_n -2 a_{n-1} +2,  a_{n-1}^2 -2 a_{n-1} +1, \\ 
        & a_{n-1} a_{n-2} -a_{n-1} -a_{n-2} +1, \\  
    & a_n a_{n-2} -2 a_{n-2},  a_{n-1} a_{n-2} -a_{n-2},  a_{n-2}^2 -a_{n-2} \rangle \enspace.
    \end{align*}
        \item Literature~\cite[p.~24, Section~2.1, Chapter~1]{Shafarevich}

        \end{itemize}
    \end{block}



\end{frame}

\begin{frame}{\GB\ diagram revisited}

\begin{figure}[h]
\centering
$$
\begin{diagram}
    \text{Data points} & \rTo^{\text{{\color<2-3>{ForestGreen} \textbf{union of single points}}}}  & \text{{\color<2-3>{ForestGreen}Ideal basis}} \\
                 & \rdTo_{\text{\MoBu}}  & \dTo_{\text{{\color<3>{dunered} \textbf{Buchberger}}}} \\
                 &  & \text{{\color<3>{dunered}\GB}} 
\end{diagram}
$$
\end{figure}

\end{frame}

\newcommand{\eigsix}{5}
\newcommand{\anqo}{7}
\newcommand{\eig}{8}

\begin{frame}{\Bua}

       From before we have: 
    \begin{align*}
    I_{12} 
    = & \langle a_n^2 -3a_n +2,  a_n a_{n-1} -a_n -a_{n-1} +1,  a_n a_{n-2} -a_n -a_{n-2} +1, \\
    & a_n a_{n-1} -a_n -2 a_{n-1} +2,  a_{n-1}^2 -2 a_{n-1} +1,  \\
        & {\color<\eigsix>{dunered} a_{n-1} a_{n-2} -a_{n-1} -a_{n-2} +1}, \\  
        & a_n a_{n-2} -2 a_{n-2}, {\color<\eigsix,\eig>{dunered} a_{n-1} a_{n-2} -a_{n-2}},  a_{n-2}^2 -a_{n-2} \rangle \enspace.
    \end{align*}
    \pause
    \begin{itemize}
        \item $(-1) \in \Q[a_n, a_{n-1}, a_{n-2}] \Rightarrow $ we can substract, \\ \pause
            \only<1-4>{i.e. $-p = (-1)\cdot p \in I_{12}$ for any $p \in I_{12}$. \\  \pause
            Follows from $0, p+q, r\cdot p \in I_{12}$ for $p,q \in I_{12}, r \in \Q[a_n, a_{n-1}, a_{n-2}]$. }  \pause
        \item {\color<\eigsix>{dunered} eighth - sixth} element of the basis:{\color<\eigsix>{dunered}  {\color<\anqo>{ForestGreen} \(q_1 = a_{n-1}-1 \in I_{12}\) } } 
            \pause
            \pause
        \item $a_n \in \Q[a_n, a_{n-1}, a_{n-2}] \Rightarrow $ we can multiply.
            \pause
        \item \(q_2 = a_n \cdot {\color<\anqo>{ForestGreen} q_1 } = a_n a_{n-1} - a_n \in I_{12}\) 
        \item \(q_3 = a_{n-2} \cdot {\color<\anqo>{ForestGreen} q_1 } = {\color<\eig>{dunered} a_{n-2} a_{n-1} - a_{n-2}} \in I_{12}\) 
            \pause
        \item
    Therefore, we can replace \(q_3\) with \(q_1\) in the basis \pause $B:= \{q_1 \text{ and the rest from the basis, except } q_3 \} $ is subset of $I_{12}$ so $ \langle B \rangle \subset I_{12}$ and 
            $q_3$ and the rest are in $\langle B \rangle$, so $I_{12} \subset \langle B \rangle $.
    \end{itemize}
    Similarly, we can further update the basis by removing the second, fourth, fifth, and sixth polynomial since they equal \(q_2 - q_1\), \(q_2 - 2q_1\), \(q_3-q_1\), and \(q_1^2\), respectively, which leads to

\end{frame}

\begin{frame}{\Bua}

    \begin{align*}
    I_{12} 
        = & \langle a_n^2 -3a_n +2,   a_n a_{n-1} -a_n   -a_{n-1} +1,  a_n a_{n-2} -a_n -a_{n-2} +1, \\
        &  a_n a_{n-1} -a_n  -2 a_{n-1} +2,  a_{n-1}^2 -2 a_{n-1} +1,  \\
        &  a_{n-1} a_{n-2}   -a_{n-1}  -a_{n-2} +1, \\  
        & a_n a_{n-2} -2 a_{n-2}, {\color{dunered} a_{n-1} a_{n-2} -a_{n-2}},  a_{n-2}^2 -a_{n-2} \rangle \enspace.
    \end{align*}

    \begin{itemize}
        \item replace {\color{dunered} \(q_3\)} with {\color{ForestGreen} \(q_1\) } in basis. \pause Similarly:
        \item second \(\mapsto {\color{blue} q_2}  {\color{ForestGreen} - q_1 }\), 
        \item fourth \(\mapsto {\color{blue}q_2}   {\color{ForestGreen} - 2 q_1 }\),
        \item fifth \(\mapsto {\color{ForestGreen} q_1^2}\), and
        \item sixth  \(\mapsto {\color{dunered}  q_3} {\color{ForestGreen} - q_1}\).
    \end{itemize}

    \begin{align*}
    I_{12} = \langle a_n^2 -3a_n +2,  a_n a_{n-2} -a_n -a_{n-2} +1, \\
    a_n a_{n-2} -2 a_{n-2},  a_{n-1} - 1,  a_{n-2}^2 -a_{n-2} \rangle \enspace.
    \end{align*}


\end{frame}

\begin{frame}{\Bua}

    \begin{align*}
    I_{12} 
        = & \langle a_n^2 -3a_n +2, {\color{blue}  a_n a_{n-1} -a_n }\ {\color{ForestGreen} - \ a_{n-1} +1},  a_n a_{n-2} -a_n -a_{n-2} +1, \\
        & {\color{blue} a_n a_{n-1} -a_n}\ {\color{ForestGreen} -\ 2 a_{n-1} +2}, {\color{ForestGreen} a_{n-1}^2 - 2 a_{n-1} +1},  \\
        & {\color{dunered} a_{n-1} a_{n-2} }\ {\color{ForestGreen} -\ a_{n-1}}\ {\color{dunered} -\ a_{n-2}}\ {\color{ForestGreen}+\ 1}, \\  
        & a_n a_{n-2} -2 a_{n-2}, {\color{dunered} a_{n-1} a_{n-2} -a_{n-2}},  a_{n-2}^2 -a_{n-2} \rangle \enspace.
    \end{align*}

    \begin{itemize}
        \item replace {\color{dunered} \(q_3\)} with {\color{ForestGreen} \(q_1\) } in basis. \pause Similarly:
        \item second \(\mapsto {\color{blue} q_2}  {\color{ForestGreen} - q_1 }\), 
        \item fourth \(\mapsto {\color{blue}q_2}   {\color{ForestGreen} - 2 q_1 }\),
        \item fifth \(\mapsto {\color{ForestGreen} q_1^2}\), and
        \item sixth  \(\mapsto {\color{dunered}  q_3} {\color{ForestGreen} - q_1}\).
    \end{itemize}

    \pause
    \begin{align*}
    I_{12} = \langle a_n^2 -3a_n +2,  a_n a_{n-2} -a_n -a_{n-2} +1, \\
        a_n a_{n-2} -2 a_{n-2}, {\color{ForestGreen} a_{n-1} - 1},  a_{n-2}^2 -a_{n-2} \rangle \enspace.
    \end{align*}

\end{frame}


\begin{frame}{\Bua}

    \begin{align*}
    I_{12} 
        = & \langle a_n^2 -3a_n +2, {\color{blue}  a_n a_{n-1} -a_n }\ {\color{ForestGreen} - \ a_{n-1} +1},  a_n a_{n-2} -a_n -a_{n-2} +1, \\
        & {\color{blue} a_n a_{n-1} -a_n}\ {\color{ForestGreen} -\ 2 a_{n-1} +2}, {\color{ForestGreen} a_{n-1}^2 - 2 a_{n-1} +1},  \\
        & {\color{dunered} a_{n-1} a_{n-2} }\ {\color{ForestGreen} -\ a_{n-1}}\ {\color{dunered} -\ a_{n-2}}\ {\color{ForestGreen}+\ 1}, \\  
        & a_n a_{n-2} -2 a_{n-2}, {\color{dunered} a_{n-1} a_{n-2} -a_{n-2}},  a_{n-2}^2 -a_{n-2} \rangle \enspace.
    \end{align*}

    Finally, we get: 

    \(I_{12} = \langle a_{n-2}^2 -a_{n-2}, {\color{ForestGreen} a_{n-1} - 1}, {\color<3>{orange} a_n -a_{n-2} -1 } \rangle\) \pause $= \langle GB \rangle $

    This is also the \GB\ of the ideal.

    \pause

    \begin{itemize}
        \item Fun fact: {\color{orange} last} - {\color{ForestGreen} $q_1$ } = $ a_n - a_{n-1} - a_{n-2} \in I_{12} $.
    \end{itemize}


\end{frame}


\begin{frame}{\Bua\ in practice}
    \begin{itemize}
        \item So far only 2 rows were used. \pause (1,1,0) and (2,1,1) 
            \pause
        \only<3>{\item our dataset:
\begin{table}[ht]
% \setlength{\tabcolsep}{7.95mm}
% \caption{Exact equation discovery data set derived from the first 8 elements of the Fibonacci sequence. \label{tab:fibomb}}
\begin{tabular}{cccc} 
\toprule
$a_n$ & $a_{n-1}$ & $a_{n-2}$
\\
\midrule
1 & 1 & 0 \\
2 & 1 & 1 \\
3 & 2 & 1 \\
5 & 3 & 2 \\
8 & 5 & 3 \\
13 & 8 & 5  \\
\bottomrule
\end{tabular}
\end{table}
            }
            \pause
        \item
             % \(\textrm{Gr\"obnerBasis}(I_{12}\cdot I_3\cdot I_4\cdot I_5\cdot I_6)\) \pause of  the product of the ideals corresponding to the first six data points results in 
             Running \Bua\ on product \(I_{12}\cdot I_3\cdot I_4\cdot I_5\cdot I_6\) \pause of ideals of six  points results in: \\
$\langle a_{n-1}a_{n-2}^2 +26a_{n-1}^2 -85a_{n-1}a_{n-2} +55a_{n-2}^2 +6a_{n-1} +29a_{n-2} -32$, \\ $  a_{n-1}^2a_{n-2} +37a_{n-1}^2 -126a_{n-1}a_{n-2} +84a_{n-2}^2 +12a_{n-1} +41a_{n-2} -49$, \\ $  a_{n-2}^3 +18a_{n-1}^2 -57a_{n-1}a_{n-2} +36a_{n-2}^2 +3a_{n-1} +20a_{n-2} -21$,  \\
            $a_n  -a_{n-1} -a_{n-2}$, \\
            $a_{n-1}^3 +51a_{n-1}^2 -183a_{n-1}a_{n-2} +126a_{n-2}^2 +23a_{n-1} +57a_{n-2} -75 \rangle$ \pause
        \item Moadeeb then filters and selects the simplest polynomials from the basis.
    \end{itemize}

\end{frame}


\begin{frame}{The end}


\begin{figure}[h]
\centering
$$
\begin{diagram}
    \text{Data points} & \rTo^{\text{{\color<2-3>{ForestGreen} \textbf{union of single points}}}}  & \text{{\color<2-3>{ForestGreen}Ideal basis}} \\
                 & \rdTo_{\text{\MoBu\cite{mb}}}  & \dTo_{\text{{\color<3>{dunered} \textbf{Buchberger\cite{bbuchphd}}}}} \\
                 &  & \text{{\color<3>{dunered}\GB}} 
\end{diagram}
$$
\end{figure}

\end{frame}

\begin{frame}[allowframebreaks]
        \frametitle{References}
        \bibliographystyle{amsalpha}
        \bibliography{moadeeb-refs.bib}


\end{frame}

% \begin{frame}
% \bibliography{moadeeb-refs}
% \end{frame}

\section{Equation discovery}


\end{document}

% % \begin{frame}{Algoritem za odkrivanje ena"cb}
% % \begin{itemize}
% % [<+->]
% %     \item dsa
% % \end{itemize}
% % \end{frame}
